% ============================================================================
% Fichier cible : ARIEL_LOGMINER_MEMOIRE_FINAL/chapters/chapitre1_introduction.tex
%   (Introduction générale, section « Positionnement et définition
%   opérationnelle des termes », PDF p. 4)
% Ligne à remplacer : 58 (« Le titre du mémoire emploie les termes
%   \textit{autonome}, \textit{distribuée} et \textit{agents intelligents
%   multitâches}. ... »).
% Auteur : Agent 1 (intégrité factuelle), chantier 2026-11-05.
% Sources : titre réel, main.tex l. 77 et 201 ; question centrale,
%   chapitre1_introduction.tex l. 95 ; classe MultiTaskIntelligentAgent,
%   chapitre6_conclusion.tex l. 97 et chapitre4_implementation.tex l. 43.
% ============================================================================

Le titre du mémoire qualifie l'architecture de \textit{multi-agents} et de \textit{légère} ; la question centrale y ajoute les qualificatifs \textit{autonome}, \textit{distribué} et \textit{multitâche}, et l'implémentation désigne ses agents comme \textit{intelligents}. Ces termes peuvent couvrir des réalités très différentes ; leur sens doit donc être précisé dans le cadre exact de ce travail.
